\section{Mapa de tipos de datos: LLM, modelos del mundo, VLA y robots}

En la entrevista se mencionan varias veces los modelos de lenguaje de gran tamaño, los modelos del mundo, VLA, el cerebro de los robots y la IA del mundo físico. Estos conceptos se confunden con facilidad. Pueden entenderse ubicándolos en un mapa de tipos de datos: los datos de texto y código son adecuados para los LLM, los datos multimodales aportan percepción, las trayectorias de robots aportan acciones, los datos de simulación aportan escenarios controlables y los pares fracaso-éxito aportan retroalimentación con alta densidad de información.

\begin{figure}[H]
\centering
\includegraphics[width=0.96\textwidth]{data-taxonomy.png}
\caption{Mapa de tipos de datos: texto, código, multimodal, trayectorias de robots, simulación y pares fracaso-éxito.}
\end{figure}

\begin{knowledgebox}{Lectura de la figura: cada modelo necesita datos distintos}
Los LLM dependen del texto y del código; los modelos del mundo necesitan datos relacionados con la predicción física; VLA necesita visión, lenguaje y acción; las políticas de robots necesitan trayectorias y retroalimentación obtenidas en el propio dispositivo. La escasez de datos no significa que no existan datos en absoluto, sino que faltan datos que correspondan a la capacidad que se busca.
\end{knowledgebox}

\subsection{Términos clave: modelos y datos}

\noindent\begin{tabularx}{\textwidth}{p{0.22\textwidth}p{0.32\textwidth}X}
\toprule
Término & Problema que resuelve & Datos que requiere \\
\midrule
LLM & Comprensión del lenguaje, razonamiento, código y tareas del mundo digital & Texto, código, corpus multimodales, datos de retroalimentación. \\
World Model & Predecir el estado del mundo físico y sus cambios futuros & Video, simulación, interacción física, datos 3D y espaciales. \\
VLA & Vision-Language-Action: conecta percepción, lenguaje y acción & Visión, instrucciones en lenguaje, trayectorias de acción, retroalimentación del entorno. \\
Physical AI & Sistemas de IA capaces de actuar en el mundo físico & Datos del cuerpo del robot, simulación, despliegue real, evaluación. \\
Behavior Benchmark & Conjunto de evaluación de largo horizonte orientado a tareas corporizadas & Entornos de simulación, definición de tareas, criterios de éxito. \\
\bottomrule
\end{tabularx}

\subsection{Relación simbiótica}

Los equipos de modelos grandes, de modelos del mundo y de VLA no están aislados. VLA puede usar un LLM base, un modelo del mundo puede funcionar como cerebro en la nube y VLA puede devolver al modelo del mundo la retroalimentación de sus acciones. Si los sistemas de evaluación convergen poco a poco, la frontera entre los modelos del mundo y VLA también podría estrecharse.

\begin{warningbox}{No confundas los modelos del mundo con VLA}
Los modelos del mundo se enfocan más en comprender y predecir el mundo físico; VLA se enfoca más en actuar en el mundo físico. Ambos coexistirán en simbiosis, pero tienen objetivos distintos, datos distintos y también formas de evaluación distintas.
\end{warningbox}

\subsection{Resumen de la sección}

El mapa de datos determina el mapa de modelos. Los LLM, los modelos del mundo, VLA y las políticas de robots necesitan tipos de datos distintos, y también se van conectando mediante la evaluación y la retroalimentación.
